% GENERATED FILE. Edit canonical lessons, then run npm run generate:books.
% canonical-source-hash: fnv1a64:8d8ce99d723cd6a4
% canonical-lessons: ZH-C233-yangzi, ZH-C233-yanjing, ZH-C233-huzi, ZH-C233-miaotiao, ZH-C233-pifu

\chapter{How People Look}
\label{ch:zh-how-people-look}
\hlchaptermodality{233}

\begin{chapteropening}
\textbf{By the end of this chapter:} \emph{I can describe someone's appearance: their glasses, a beard, a slim figure and their skin.}

Complete the last lesson of chapter 233: I can describe someone's appearance: their glasses, a beard, a slim figure and their skin.
\end{chapteropening}

\section[yàngzi]{yàngzi — appearance; manner}

\label{lesson:ZH-C233-yangzi}



\begin{quote}
Before the new one: say the Mandarin for to crash into, then the Mandarin for an ambulance.
\end{quote}

\subsection*{You'll want to know: yàngzi}
\emph{yàngzi} — \textquotedblleft{}appearance; manner\textquotedblright{}.

\begin{quote}
\emph{tā de yàngzi hěn xiàng tā bàba} — He takes after his dad.
\end{quote}

Its characters come later; for now it is a word for the ear.

\begin{cousinweb}
\emph{yàngzi} is built from \emph{yàng} 'shape' and the noun ending \zh{子}, inherited from Middle Chinese \emph{tsiX} (Guangyun, in Baxter's transcription). No dictionary we use records an older form of the first, so its history is not settled.
\end{cousinweb}

\begin{grammarlens}[title={What you've built}]
One more word for describing things.
\end{grammarlens}

\subsection*{Your turn}
\begin{itemize}
\raggedright
  \item \emph{Say it:} \emph{yàngzi}
  \item \emph{Say it:} \emph{yàngzi}, once more
  \item \emph{Say it:} \emph{tā de yàngzi hěn xiàng tā bàba} — He takes after his dad.
  \item \emph{Recall:} say the Mandarin for to crash into, then the Mandarin for an ambulance, then say \emph{yàngzi} again
\end{itemize}

\begin{tcolorbox}[breakable,colback=teal!4,colframe=teal!35!black,title={Before you move on}]
What does yàngzi mean? (\textquotedblleft{}Appearance; manner\textquotedblright{}.) Say it once more.
\end{tcolorbox}
\section[yǎnjìng]{yǎnjìng — glasses; spectacles}

\label{lesson:ZH-C233-yanjing}



\begin{quote}
Before the new one: say the Mandarin for an ambulance, then the Mandarin for appearance.
\end{quote}

\subsection*{You'll want to know: yǎnjìng}
\emph{yǎnjìng} — \textquotedblleft{}glasses; spectacles\textquotedblright{}.

\begin{quote}
\emph{wǒ de yǎnjìng ne} — Where are my glasses?
\end{quote}

Its characters come later; for now it is a word for the ear.

\begin{cousinweb}
\emph{yǎnjìng} is built from \emph{yǎn} 'eye' and \emph{jìng} 'mirror', both inherited from Middle Chinese, \emph{ngeanX} and \emph{kjaengH} (Guangyun, in Baxter's transcription), which glosses the first 'the eye'.
\end{cousinweb}

\begin{grammarlens}[title={What you've built}]
One more word for describing things.
\end{grammarlens}

\subsection*{Your turn}
\begin{itemize}
\raggedright
  \item \emph{Say it:} \emph{yǎnjìng}
  \item \emph{Say it:} \emph{yǎnjìng}, once more
  \item \emph{Say it:} \emph{wǒ de yǎnjìng ne} — Where are my glasses?
  \item \emph{Recall:} say the Mandarin for an ambulance, then the Mandarin for appearance, then say \emph{yǎnjìng} again
\end{itemize}

\begin{tcolorbox}[breakable,colback=teal!4,colframe=teal!35!black,title={Before you move on}]
What does yǎnjìng mean? (\textquotedblleft{}Glasses; spectacles\textquotedblright{}.) Say it once more.
\end{tcolorbox}
\section[húzi]{húzi — beard; moustache}

\label{lesson:ZH-C233-huzi}



\begin{quote}
Before the new one: say the Mandarin for appearance, then the Mandarin for glasses for the eyes.
\end{quote}

\subsection*{You'll want to know: húzi}
\emph{húzi} — \textquotedblleft{}beard; moustache\textquotedblright{}.

\begin{quote}
\emph{wǒ bàba yǒu húzi} — My dad has a beard.
\end{quote}

Its characters come later; for now it is a word for the ear.

\begin{cousinweb}
\emph{húzi} is built from \emph{hú} 'beard' and the noun ending \zh{子}, inherited from Middle Chinese \emph{tsiX} (Guangyun, in Baxter's transcription). No dictionary we use records an older form of the first, so its history is not settled.
\end{cousinweb}

\begin{grammarlens}[title={What you've built}]
One more word for describing things.
\end{grammarlens}

\subsection*{Your turn}
\begin{itemize}
\raggedright
  \item \emph{Say it:} \emph{húzi}
  \item \emph{Say it:} \emph{húzi}, once more
  \item \emph{Say it:} \emph{wǒ bàba yǒu húzi} — My dad has a beard.
  \item \emph{Recall:} say the Mandarin for appearance, then the Mandarin for glasses for the eyes, then say \emph{húzi} again
\end{itemize}

\begin{tcolorbox}[breakable,colback=teal!4,colframe=teal!35!black,title={Before you move on}]
What does húzi mean? (\textquotedblleft{}Beard; moustache\textquotedblright{}.) Say it once more.
\end{tcolorbox}
\section[miáotiao]{miáotiao — slim; slender}

\label{lesson:ZH-C233-miaotiao}



\begin{quote}
Before the new one: say the Mandarin for glasses for the eyes, then the Mandarin for a beard.
\end{quote}

\subsection*{You'll want to know: miáotiao}
\emph{miáotiao} — \textquotedblleft{}slim; slender\textquotedblright{}.

\begin{quote}
\emph{tā hěn miáotiao} — She's very slender.
\end{quote}

Its characters come later; for now it is a word for the ear.

\begin{cousinweb}
\emph{miáotiao} is built from \emph{miáo} 'a sprout' and \emph{tiáo} 'a twig'. The second is inherited from Middle Chinese \emph{dew} (Guangyun, in Baxter's transcription), glossed 'a small branch'; the dictionary glosses the first, \emph{mjew} there, 'young crops in the field'. In the word the second loses its tone.
\end{cousinweb}

\begin{grammarlens}[title={What you've built}]
One more word for describing things.
\end{grammarlens}

\subsection*{Your turn}
\begin{itemize}
\raggedright
  \item \emph{Say it:} \emph{miáotiao}
  \item \emph{Say it:} \emph{miáotiao}, once more
  \item \emph{Say it:} \emph{tā hěn miáotiao} — She's very slender.
  \item \emph{Recall:} say the Mandarin for glasses for the eyes, then the Mandarin for a beard, then say \emph{miáotiao} again
\end{itemize}

\begin{tcolorbox}[breakable,colback=teal!4,colframe=teal!35!black,title={Before you move on}]
What does miáotiao mean? (\textquotedblleft{}Slim; slender\textquotedblright{}.) Say it once more.
\end{tcolorbox}
\section[pífū]{pífū — skin}

\label{lesson:ZH-C233-pifu}



\begin{quote}
Before the new one: say the Mandarin for a beard, then the Mandarin for slender.
\end{quote}

\subsection*{You'll want to know: pífū}
\emph{pífū} — \textquotedblleft{}skin\textquotedblright{}.

\begin{quote}
\emph{tā de pífū hěn bái} — She has very fair skin.
\end{quote}

Its characters come later; for now it is a word for the ear.

\begin{cousinweb}
\emph{pífū} is built from \emph{pí} 'hide' and \emph{fū} 'skin'. The first is inherited from Middle Chinese \emph{bje} (Guangyun, in Baxter's transcription), glossed with this very word, 'skin'; the dictionary glosses the second, \emph{pju} there, the same way.
\end{cousinweb}

\begin{grammarlens}[title={What you've built}]
One more word for describing things.
\end{grammarlens}

\subsection*{Your turn}
\begin{itemize}
\raggedright
  \item \emph{Say it:} \emph{pífū}
  \item \emph{Say it:} \emph{pífū}, once more
  \item \emph{Say it:} \emph{tā de pífū hěn bái} — She has very fair skin.
  \item \emph{Recall:} say the Mandarin for a beard, then the Mandarin for slender, then say \emph{pífū} again
\end{itemize}

\begin{tcolorbox}[breakable,colback=teal!4,colframe=teal!35!black,title={Before you move on}]
What does pífū mean? (\textquotedblleft{}Skin\textquotedblright{}.) Say it once more.
\end{tcolorbox}
